%! Author = lili
%! Date = 7/9/26

\chapter{Tutorium 08.06.2026}
\untit{Präsenzaufgaben 2.1: Erwartungswert und Varianz \hfill 8.--12.\ Juni 2026}
\txo{Bearbeiten Sie wahlweise Aufgabe 2.1.4 oder 2.1.5. Wenn Zeit bleibt, bearbeiten Sie
zusätzlich Aufgabe 2.1.6.}

\section{PÜ 2.1.4}
Sei $X$ eine Zufallsvariable mit Bildbereich $X(\eraum) = \{-2,-1,0,1,2\}$, die auf
einem Wahrscheinlichkeitsraum $(\eraum,\wfk)$ definiert ist. Es gelte
\[
    \pe(X = -1) = \pe(X = 1) = \frac{p}{4}
    \qquad\text{und}\qquad
    \pe(X = -2) = \pe(X = 2) = \frac{1}{4}
\]
mit einem Parameter $p$.

\begin{auf}[Bestimmen Sie die Verteilung $\wfk_X$ von $X$ in Abhängigkeit von $p$. Welche
Werte von $p$ sind zulässig?]
    NOCH KEINE LÖSUNG % TODO NOCH KEINE LÖSUNG
\end{auf}

\begin{auf}[Existiert der Erwartungswert von $X$? Falls ja, berechnen Sie diesen.]
    NOCH KEINE LÖSUNG % TODO NOCH KEINE LÖSUNG
\end{auf}

\begin{auf}[Existiert die Varianz von $X$? Falls ja, berechnen Sie diese.]
    NOCH KEINE LÖSUNG % TODO NOCH KEINE LÖSUNG
\end{auf}

\begin{auf}[Seien nun $p = \frac{1}{2}$ und $Y = -2X + 4$. Bearbeiten Sie die
Aufgabenteile (b) und (c) auch für $Y$.]
    NOCH KEINE LÖSUNG % TODO NOCH KEINE LÖSUNG
\end{auf}

\section{PÜ 2.1.5}
Die gemeinsame Verteilung zweier Zufallsvariablen $X$ und $Y$ sei gegeben durch
\[
    \begin{array}{c|ccc}
        \wfk_{(X,Y)}\left((x,y)\right) & y = -2 & y = 0 & y = 2 \\ \hline
        x = -1 & \frac{1}{8} & c & 2c \\[2pt]
        x = 1  & \frac{1}{4} & 0 & \frac{1}{4}
    \end{array}
\]
Außerdem sei $Z = 2X + Y$.

\begin{auf}[Bestimmen Sie $c$.]
    NOCH KEINE LÖSUNG % TODO NOCH KEINE LÖSUNG
\end{auf}

\begin{auf}[Geben Sie die Bildbereiche $X(\eraum)$, $Y(\eraum)$ und $Z(\eraum)$ an.]
    NOCH KEINE LÖSUNG % TODO NOCH KEINE LÖSUNG
\end{auf}

\begin{auf}[Begründen Sie, warum Erwartungswert und Varianz von $X$, von $Y$ und von $Z$
existieren.]
    NOCH KEINE LÖSUNG % TODO NOCH KEINE LÖSUNG
\end{auf}

\begin{auf}[Berechnen Sie die Randverteilungen $\wfk_X$ von $X$ und $\wfk_Y$ von $Y$.]
    NOCH KEINE LÖSUNG % TODO NOCH KEINE LÖSUNG
\end{auf}

\begin{auf}[Berechnen Sie Erwartungswert und Varianz von $X$ und von $Y$.]
    NOCH KEINE LÖSUNG % TODO NOCH KEINE LÖSUNG
\end{auf}

\begin{auf}[Berechnen Sie den Erwartungswert von $Z$.]
    NOCH KEINE LÖSUNG % TODO NOCH KEINE LÖSUNG
\end{auf}

\begin{auf}[Berechnen Sie die Wahrscheinlichkeit, dass $Z > 0$ gilt.]
    NOCH KEINE LÖSUNG % TODO NOCH KEINE LÖSUNG
\end{auf}


\section{PÜ 2.1.6}
In Fortsetzung von Aufgabe 1.2.4 betrachten wir noch einmal die Dichte
\[
    f(x) =
    \begin{cases}
        3c\,x^2(1 - x) & \text{für } x \in [0,1],\\[2pt]
        0 & \text{sonst,}
    \end{cases}
\]
wobei $c \in \rz$ die in der zweiten Präsenzübung bestimmte Konstante sei. Sei $X$
eine Zufallsvariable mit Dichte $f$.

\begin{auf}[Berechnen Sie Erwartungswert und Varianz von $X$. Warum existieren diese?]
    NOCH KEINE LÖSUNG % TODO NOCH KEINE LÖSUNG
\end{auf}

\begin{auf}[Berechnen Sie Erwartungswert und Varianz von $-4X - 1$. Warum existieren diese?]
    NOCH KEINE LÖSUNG % TODO NOCH KEINE LÖSUNG
\end{auf}
% Blatt

% TODO Übung 2.1.1-2.1.3 bilder nils